% Equazione di continuità: tubo che si restringe. Nello stesso tempo passa lo stesso
% volume: cilindro corto e largo nella sezione S1, lungo e stretto nella sezione S2.
\begin{tikzpicture}[>=Stealth, font=\small]
  % tubo
  \fill[formula!10] (0,-0.8) -- (2.4,-0.8) .. controls (3.0,-0.8) and (3.0,-0.35) .. (3.6,-0.35)
    -- (8.2,-0.35) -- (8.2,0.35) -- (3.6,0.35) .. controls (3.0,0.35) and (3.0,0.8) .. (2.4,0.8)
    -- (0,0.8) -- cycle;
  % volumi uguali
  \fill[esempio!40] (0.7,-0.8) rectangle (1.3,0.8);
  \fill[esempio!40] (4.3,-0.35) rectangle (7.4,0.35);
  \draw[esempio, thick] (0.7,0) ellipse (0.15 and 0.8);
  \draw[esempio, thick] (4.3,0) ellipse (0.07 and 0.35);
  % pareti
  \draw[very thick] (0,-0.8) -- (2.4,-0.8) .. controls (3.0,-0.8) and (3.0,-0.35) .. (3.6,-0.35) -- (8.2,-0.35);
  \draw[very thick] (0,0.8) -- (2.4,0.8) .. controls (3.0,0.8) and (3.0,0.35) .. (3.6,0.35) -- (8.2,0.35);
  % velocità
  \draw[->, very thick, formula] (1.5,0) -- (2.2,0) node[right] {$v_1$};
  \draw[->, very thick, formula] (7.5,0) -- (9.4,0) node[right] {$v_2$};
  % etichette
  \node[esempio] at (0.7,1.1) {$S_1$};
  \node[esempio] at (4.3,0.65) {$S_2$};
  \draw[|<->|] (0.7,-1.1) -- (1.3,-1.1) node[midway, below] {$v_1 \Delta t$};
  \draw[|<->|] (4.3,-0.65) -- (7.4,-0.65) node[midway, below] {$v_2 \Delta t$};
\end{tikzpicture}
